\section{Discussion}
\label{sec:discussion}

\subsection{What Including Debate Reveals About Role-Playing Games}
\label{sec:rpg-insights}

The meeting of playtest and theoretical expectations gave birth to a handful of insights:

Firstly, on alibi. The framework currently treats alibi as a single concept appearing under two complementary glosses. \textcite{deterding2017alibis} names the \textit{justification-alibi}: the rationale an adult adopts as licence to play in the first place. \textcite{bowman2025transformative} names the \textit{performance-alibi}: the in-play permission to act, decide, and feel under cover of character. The Inside Out playtests show these are not parallel but coupled. The Rolling Out closing ritual, a framing device of the kind most LARPs use without difficulty, was reported as ``robotic'' in the first session and met with mumbling in the second. Yet the same participants spontaneously adopted ``Inner Voices'' as in-play address. Both moves were formally licensed by the same in-play permission. What discriminated them was justification: chanting felt decorative under a rationale that read play as ``learning argumentation,'' but natural where the rationale was ``inhabiting a character.'' The justification an adult brings to the table places a prior on the performance-alibis they will exercise once seated. The contribution to the framework is to treat the adult-justification surface not as a separate designable surface from the in-play one, but as a filter on it.

The second refinement concerns evaluator roles. Reform discourse around debate's value-convergence problem has long treated judges as neutral arbiters in need of better calibration. The role-playing-game framing identifies a different mechanism. Standard formats ask every judge to roleplay the same archetype, the average, unbiased person, and that is what produces drift. Conflating unbiased with lacking values procedurally teaches that certain values are a baseline and thus pulls people to adhere to them. Inside Out's playtests proved that its structural choices can directly prevent that. By prescribing for the Subject a clearly biased character and asking players to also voice their interpretative biases over that character's beliefs during the S1 speech, it led to players' realization of the need for tolerance towards others' beliefs and circumstantial decisions. This becomes thus a clean demonstration of the thesis's methodological claim: only treating the judge as a roleplayed character names the parameter that needed redesign, where the standard reform vocabulary of bias and calibration does not. 

This is further substantiated by the manner in which the apparent tension in the Subject role (emotional investment and evaluative neutrality) was proved to be easier resolved by leaning into roleplay in a \textit{steered} \parencite{montola2015steering, pohjola2015steering} fashion. This is a mirror of my and Seo's (2024) thesis that the cognitive overload of debate is only manageable by genuine character inhabiting. This process of steered roleplay proves thus to be employable in the pursuit of performance and with it towards memetic bleed, in itself a teachable capacity that can help teach other capacities (the long term cognitive abilities that form the transformative potential of debate).

The fourth refinement concerns where pressure comes from once external competition has been removed. The framework anticipated that removing win/loss outcomes would require a replacement source of competitive motivation, and the commendation architecture was designed to provide it. The playtests instead surfaced alongside them the alternative mechanism of social recognition one's performance. This mirrors the type of performance pressure found in LARP communities in which there is no formal performance metric. This raises questions on whether container safety is better reached by enforcing a culture of fake appraisals and hushed judgements, or one of open critique of players' performance, shaped constructively now that it is in the open.

All these lead not just to the pedantic overproof of debate as a roleplaying game. Debate is a special category of roleplaying games - it is evaluative, unavoidably somewhat conflictual and naturally more performance demanding than many roleplaying games. By stopping debating it is that and focusing on theorizing its inner workings, we discover further insights about other roleplaying games as well.

\subsection{Directions for Continued Study}
\label{sec:further-study}

The playtests open four research directions, each pointing to an established literature outside transformative-game-design ready to be imported. First, although the commendation system was received enthusiastically, evaluators struggled to assign commendations consistently and named a risk of ``farming'': stylised play optimised for commendation coverage rather than development. The framework specifies what transformation \textit{is} but neither how a facilitator sees it in real time nor how a reward structure can be designed so that optimising for it is identical to the development it registers, an alignment problem on which gamification and progressive-game-design literatures already have substantial theory. 

Second, participants' reconstruction of how speedtalking became British Parliamentary's meta and contaminated World Schools located the cause not in rules but in who held social power in the circuit once the founding teachers disengaged; governance, how rules are interpreted, evaluators trained, motions authored, and meta drift sustained, is a designable surface above the format, where community-management literature offers tools that \textcite{bowman2025transformative}'s ``community support'' category does not yet capture. 

Third, naming peer recognition as a pressure category opens questions the playtests cannot answer: how it scales beyond the bounded session, how it interacts with progression, and what failure modes (clique formation, in-group bias, the homogenising drift a unitary evaluator role produces in standard debate) it carries.

Fourth, Inside Out's design moves are decomposable: the commendation system, the split between an evaluated Interlocutor and an unevaluated external Evaluator, three teams, a judge who roleplays a named character, and teams prescribed values but granted creative freedom over their expression are each candidates for single-mechanism transplantation, inviting a methodology that treats transformative mechanisms as transplantable units rather than constituents of holistic experiences. 

None of these gaps is debate-specific; they apply to any transformative format that removes external competition, sustains itself through a community, and must register development without distorting it. Importing progression theory, community management, peer-recognition dynamics, and mechanic-level analysis is the research programme this thesis sharpens but that the field would benefit from regardless of Inside Out's continued iteration.

Together with specific mechanical refinement, all these coalesce into a research programme that the next iteration of work, on Inside Out or elsewhere, can now pursue.
